\scp{Unclassified languages}{07-05}
Before we discuss each of the subgroups within central
Maluku in more detail,
we need to draw attention to two languages which
cannot yet be classified on the basis of shared phonological
innovations due to lack of data.
These are Salas, of eastern Seram and Hukumina of northwest Buru.

\sscp{Salas [sgu]}{07-05-01}
According to \citet[121f]{LoskiLoski1989} “Salas Gunung is
spoken by 50 people in the village by that name in the Waru Bay.
Most of the people are reported to be bilingual in the Masiwang language.
\citet{Tauern1931} presents a short grammar of Liambata,
possibly a dialect of Salas Gunung”.
Salas is placed by \citet{LoskiLoski1989} within their “Manusela-Seti
language chain” and thus within our \tsc{Ambon-Seram},
and not with Bobot or Masiwang,
which belong in our \tsc{Seram-Tanimbar-Bomberai}.

Our examination of the ``Liambata'' data in \citet[128--139]{Tauern1931}
indicates that it may be a variety of Seti,
based on distinctive lexical items such as \it{inta} ‘skin’ (Liana-Seti
\it{inta} ‘skin’), and \it{ufeia} ‘house’ (Liana-Seti \it{ufai-a} ‘house’).
At the same time,
this data diverges from other sources of Seti in several respects
indicating that it may indeed be a different language.
It is important to note that the data in \citet[128--139]{Tauern1931}
is different from the “Liambata” data in \citet{Stresemann1927},
the latter of which \it{is} (almost certainly) a variety of Seti.

Given the equivocation of \citet{LoskiLoski1989} regarding the
data in \citet{Tauern1931}, as well as our own judgment that the data from \citet{Tauern1931}
is related to but different from Seti,
we cannot confidently identify any Salas data to make an assessment of its position.
Lexically, it is closer to \tsc{Ambon-Seram} languages than \tsc{Seram-Tanimbar-Bomberai}.

\sscp{Hukumina [huw]}{07-05-02}
Hukumina is an extinct language of north-western Buru.
\citet[63--76]{Stokhof1982b} contains a word list labelled ``Hukumina'',
but the reliability of this list is very low,
and it contains many items which are clearly Buru and Kayeli.
Grimes was able to collect a small amount of data in 1989
and states:

\begin{quote}
	``Hukumina was referred to by \citet{Hendriks1897} as a dialect of Buru spoken in the districts
	of Hukumina, Palumata and Tomahu (these were districts in the north-west comer of Buru).
	He gives no other details other than his impression that one could shift easily between
	Hukumina and other dialects of Buru. A Hukumina word list is published in the Holle lists
	(Stokhof 1982) [\cite{Stokhof1982b}] with very little information other than the date 1896.
	
	[\ldots] I was able to find one old woman from the former village of Hukumina that used to be
	located behind the present village of Masarete near the fort at Kayeli (north-east Buru). My
	data are even more divergent than the Holle list (the latter includes several items that are
	clear borrowings from Buru). The language of Hukumina is very distinct from Buru in
	intonation, lexicon and sound correspondences. For example, PAn *k is Buru /k/ and
	Hukumina /c/ intervocalically.'' \citep[99]{Grimes2000a}
\end{quote}

\begin{table}
	\centering
	\caption[Buru, Kayeli, and Hukumina words (Grimes 2010: 78)]
	{Buru, Kayeli, and Hukumina words \citep[78]{Grimes2010b}}
	\label{tab:07-19}
	\st{6pt}\begin{tabular}{llll}\lsptoprule
Gloss	&	Buru	&	Kayeli	&	Hukumina	\\	\midrule
hand, arm	&	fahan	&	limani	&	rima	\\	
mouth	&	fifin	&	nuan	&	mpidu-	\\	
foot, leg	&	kadan	&	bitiɲ	&	eɲaik	\\	
head	&	olon	&	oloni	&	fatun	\\	
betelpepper	&	dalu	&	gamal	&	elut	\\	
jungle	&	mua	&	boot	&	gulalin	\\	
sleep	&	bage	&	ine	&	etnono	\\	
sleepy	&	duba	&		&	lastoton	\\	
sick, painful	&	empei	&	berere	&	boto	\\	
return (home)	&	oli	&	wae	&	bihin	\\	
no, not	&	moo	&		&	mbaa	\\	
hungry	&	eglaba	&	nanibar	&	spala	\\	
sit	&	defo	&	stea	&	egnabon	\\	
		\lspbottomrule
	\end{tabular}
\end{table}

He further states that the consultant's ``[\ldots] mind wandered
regularly, and the little data I was able to collect from her are a mixture of Kayeli,
Buru, and some other language that I assume is Hukumina.'' \citep[75]{Grimes2010}
A sample of Hukumina vocabulary compared with Buru and Kayeli is given in \trf{07-19}
to illustrate the difference between this language and others of Buru Island.

Grimes' suspicion is that Hukumina may have been a long-standing colony from
Southeast Sulawesi, as there are thousands of migrants from there along the
west and south coasts of Buru. However, the lack of reliable data means
that this cannot be confirmed. Unless more data becomes available,
Hukumina remains unclassifiable. We list it with other \tsc{Sula-Buru} languages
in \crf{ch:SB} and appendix \ref{ch:LanLis} out of convenience according to its geographic location.
